% Generated draft: requires mathematical and source review.
\documentclass[11pt]{article}
\usepackage[margin=25mm]{geometry}
\usepackage{fontspec}
\setmainfont{Latin Modern Roman}
\usepackage{amsmath,amssymb,mathtools,unicode-math}
\setmathfont{Latin Modern Math}
\usepackage{enumitem,needspace,xurl,hyperref}
\setlength{\parindent}{0pt}
\setlength{\parskip}{5pt}
\setlength{\emergencystretch}{3em}
\newcommand{\N}{\mathbb{N}}
\newcommand{\Z}{\mathbb{Z}}
\newcommand{\Q}{\mathbb{Q}}
\newcommand{\R}{\mathbb{R}}
\newcommand{\C}{\mathbb{C}}
\newcommand{\F}{\mathbb{F}}
\newcommand{\A}{\mathbb{A}}
\newcommand{\PP}{\mathbb{P}}
\newcommand{\E}{\mathbb{E}}
\newcommand{\eps}{\varepsilon}
\newcommand{\dd}{\,\mathrm{d}}
\begin{document}
\section*{UnsolvedMath ID: 30004175}
\textbf{Source problem code:} catalogue entry
\subsection{Definitions and mathematical statement}
Let \((X,J)\) be a complex manifold of complex dimension \(n\). A Kähler metric \(g\) has an associated real, closed, positive \((1,1)\)-form \(\omega\), given by \(\omega(u,v)=g(Ju,v)\). It is scalar-flat if its scalar curvature satisfies \(\operatorname{Scal}(g)\equiv 0\). Write \([\omega]_{\mathrm{dR}}\in H^2_{\mathrm{dR}}(X;\R)\) for its ordinary de Rham cohomology class.

The abbreviation ALE means asymptotically locally Euclidean: outside a compact set, the metric is asymptotic to a Euclidean quotient by a finite group acting freely on the sphere at infinity. The supplied record does not specify the dimension, completeness requirement, end conditions, decay rate, derivative estimates, asymptotic coordinates, or normalization at infinity. These conditions must be fixed before the admissible class of metrics is fully determined.

\textbf{Uniqueness question, with unresolved qualifications.} For a fixed ALE Kähler complex manifold \((X,J)\) admitting scalar-flat Kähler metrics, and a Kähler class \(\alpha\) containing such a metric, is the scalar-flat Kähler metric in \(\alpha\) unique? Under the literal interpretation of uniqueness as equality of metrics on the fixed complex manifold, and of Kähler class as ordinary de Rham class, the proposed implication is
\[
\left.
\begin{gathered}
 g_0,g_1\text{ are admissible scalar-flat ALE Kähler metrics on }(X,J),\\
 [\omega_0]_{\mathrm{dR}}=[\omega_1]_{\mathrm{dR}}=\alpha
\end{gathered}
\right\}
\quad\Longrightarrow\quad g_0=g_1\;?
\]
Here admissibility refers to the still unspecified asymptotic and completeness conditions. The source does not determine whether uniqueness instead means uniqueness modulo a specified group of holomorphic automorphisms, or whether the class includes additional asymptotic data. Thus this implication records a literal reading of the question, rather than a fully resolved formulation. No existence assertion is made for classes that do not contain a scalar-flat Kähler metric.
\subsection{Short English statement}
On a fixed ALE Kähler manifold, must two scalar-flat Kähler metrics in the same Kähler class coincide? The supplied record leaves the asymptotic conditions, the meaning of Kähler class on this noncompact manifold, and possible equivalence under automorphisms unspecified. Its open label is supplied curation, not an independently verified status.
\subsection{Sources}
\begin{itemize}
\item Supplied local registry record: permanent UnsolvedMath ID \(30004175\), title \emph{Uniqueness of Scalar-Flat Kähler ALE Metrics}, exact source problem code \texttt{catalogue entry}, in the Oberwolfach Reports Open Problems collection. Project registry references are \texttt{data/unsolvedmath/problems.json} and \texttt{catalog/problems.json}; their file contents were not independently inspected for this draft.
\item Supplied report citation: \emph{Differentialgeometrie im Großen (2020)}, Oberwolfach Reports, DOI \texttt{10.4171/owr/2019/30}, pp.\ 1791--1839. Supplied source URL and literature-source link, labelled ``original OWR report'': \url{https://doi.org/10.4171/owr/2019/30}. The record describes the collection as drawn from workshop open-problem discussions. No individual proposer is supplied; the proposed year is recorded as 2019.
\item Supplied literature assessment dated \texttt{2026-08-22}: status ``open.'' The assessment says that the record contains only an assertion that a problem is completely open and omits its mathematical statement, so no stronger classification is justified. This conflicts with the explicit uniqueness question present in the supplied record. The date and assessment are supplied dated curation, not independently verified literature findings. The background expressly requests independent verification of time-sensitive status and cited claims.
\item The supplied statement status is \texttt{corrected\_verified}; its verification note reports manual semantic, typographic, and LaTeX review confirming a self-contained statement preserving the target without unsupported additions. This is a supplied review claim and does not resolve the formulation gaps identified here.
\end{itemize}
\end{document}
